\begin{afterword}
\summaryhead

The post opens with a child asking where a rock came from, then the boulder, the mountain, the giant Ymir and the abyss Ginnungagap, until the answer is ``Never ask that question.'' It then turns to the puzzle of the First Cause. Science traces events back to the Big Bang, but the physical laws themselves still seem to call for an explanation. Some people answer ``God!'' The author answers that this does not help, since one can ask where God came from, and ``God is uncaused'' leaves the same question as ``Time began with the Big Bang.''

Following Jonathan Wallace, the author suggests that ``God!'' works as a semantic stopsign: not a claim, but a signal not to think past this point. Other words do the same for other people. An acquaintance of the author answers every policy question about new technologies with ``Liberal democracy!'', though he would question any other answer. The post ends with a caution: no word is a stopsign in itself, and what marks one is a person's failure to consider the obvious next question.

\responsehead

The post names an idea and credits it. The name is Jonathan Wallace's, and the post's footnote gives Wallace's essays. Its last paragraph states a fair test, whether a person asks of their own answer the questions they would ask of any other, and warns against using the idea as a ready-made reply to views one dislikes. The problems are in its examples.

The lead example leaves out the answer its targets give. The post asks what could make anyone think ``God!'' helps with the First Cause, and replies to ``God is uncaused'' that we would then ask ``why some events but not others are allowed to be uncaused.'' Defenders of the cosmological argument answer that question. On the argument from contingency, only what is contingent needs an explanation; on the kalam version, only what begins to exist needs a cause. The reply is disputed, but the post does not mention it. It explains the answer ``God!'' instead by a motive, that saying it ``is a way of belonging to a tribe,'' and offers no evidence for that.

By the post's own test, the theists of its opening example are not using a stopsign. A theist who answers ``Where did God come from?'' with the distinction between contingent and necessary things has considered the obvious next question, whether or not the answer is sound. By the definition in the last paragraph, that person is not using a stopsign. The post's criterion is more careful than its first example.

The modern example cannot be checked. It is one unnamed acquaintance, known only through the author's summary, and the second half of the test is imagined: the post reports what he ``would have asked'' of another answer, not anything he said. The remaining cases are hypotheticals. The reader is given a test, but no case in which it was applied to what a named person actually said.

In short: a credited idea with a fair test at the end; by that test, the theists who give the standard reply in the post's lead example are not using a stopsign, and the modern example is one unnamed acquaintance and an imagined reply.
\end{afterword}
